% ===================================================================
\section{Eutrophication and alternative stable states (EE06)}
% ===================================================================

\subsection{Lake types: boxes and a continuum}
\textbf{Classical} (Naumann, Thienemann): \textbf{oligotrophic} (low productivity, clear water), \textbf{eutrophic}
(high productivity, green water from phytoplankton), \textbf{dystrophic} (low conductivity, brown humic water), then
``extinction'' as the basin fills with sediment. Qualitative signs: oligotrophic lakes have an \emph{orthograde} summer
oxygen profile (oxygen stays high with depth), eutrophic lakes a \emph{clinograde} one (oxygen falls to zero in the
hypolimnion).

\textbf{Modern}: oligo- and eutrophic are regions of a \textbf{continuum}, separated by limits such as $P_{tot}$ during
spring circulation (µg\,L$^{-1}$): ultra-oligotrophic $<5$, oligotrophic 5--10, mesotrophic 10--30, eutrophic 30--100,
hypereutrophic $>100$. Since trophic state is a biological phenomenon, several criteria are needed.

\begin{trap}[: one number can misclassify]
Schöhsee (mean depth 13\,m, $P_{tot}\approx20$) is mesotrophic by phosphorus but has an anoxic hypolimnion
(``eutrophic'' traditionally). Lake Constance (mean depth $\sim$100\,m, $P_{tot}\approx100$ in 1980) is eutrophic by
phosphorus yet its bottom water stayed $\sim$50\,\% saturated. The difference is \textbf{morphometry}: a huge hypolimnion
holds enough oxygen from spring overturn to decompose sinking matter; a small one does not.
\end{trap}

\textbf{What sets trophic status} (slide 5): (1) basin shape (mean depth); (2) the drainage basin (erosion, soils);
(3) climate (precipitation, temperature); (4) human activity (\textbf{cultural eutrophication}: fertiliser, sewage,
industrial waste). Large deep lakes are less productive because (i) loading per unit volume is lower when the lake is
big relative to its watershed, (ii) little of the epilimnion touches the sediment, so few nutrients return to the lit
zone, and (iii) longer retention times let more phosphorus settle out. Background (OECD/Vollenweider):
$P_L=1.55[P_{in}/(1+\sqrt{\tau_w})]^{0.82}$.

\subsection{How a lake becomes eutrophic}
\begin{chain}[: the eutrophication chain, and why it self-accelerates]
\begin{enumerate}
\item \textbf{Nutrient enrichment}, chiefly phosphorus (sewage, P-detergents, fertiliser, erosion). P limits because it
is relatively immobile in soil, unlike N. Sewage has N:P $\approx$ 4:1, so heavy sewage can even shift a lake to
N limitation.
\item \textbf{More production}: algal blooms, including cyanobacteria (\emph{Anabaena}, \emph{Microcystis}): scums,
toxins, odours. Periphyton overgrows and shades submerged plants.
\item \textbf{Less light}; dead algae sink and decompose, consuming \textbf{oxygen in the hypolimnion}.
\item \textbf{Anoxic sediment releases phosphorus} (internal loading; background chemistry: Fe(III) phosphate is
reduced to soluble Fe(II) and phosphate).
\item More P $\Rightarrow$ more algae $\Rightarrow$ more anoxia $\Rightarrow$ more P: every link positive, so the loop is
\textbf{reinforcing}. Internal loading ``sets in motion the self-acceleration of eutrophication''.
\item pH rises (intense photosynthesis), turning ammonium into toxic ammonia; oxygen collapses: \textbf{fish kills}.
\end{enumerate}
Not only a small-lake problem: Caspian Sea (2003), Lake Erie. Anoxic hypolimnia exist only in stratified lakes, which is
why a deep lake with a big hypolimnion tolerates more phosphorus.
\end{chain}

\subsection{Reversal: Lake Washington and the lag}
Lake Washington (Seattle): sewage from treatment plants in the 1940s--50s; in 1955 \emph{Oscillatoria rubescens}
appears (Edmondson's alarm); early 1960s transparency $\sim$1\,m (``Lake Stinko''); sewage diversion to Puget Sound
complete \textbf{1968}; phosphorus falls within $\sim$2 years; transparency 3--4\,m (1970--75); \textbf{1976}: large
\emph{Daphnia} return and transparency doubles to $\sim$10\,m. (Slide 11's ``sewage $\to$ oligotrophy $\to$ Lake
Stinko'' is a slip: sewage made it eutrophic.) Lake Constance went from 6 to 87 to 8\,µg\,L$^{-1}$ with tertiary
treatment and P-free detergents.

\begin{chain}[: why the second jump was not phosphorus]
Phosphorus had already fallen by 1970, giving 3--4\,m. The jump to 10\,m came with the return of large \emph{Daphnia},
which graze phytoplankton. They returned because filamentous cyanobacteria that clog their filters disappeared
(bottom-up) and mysid shrimp predation fell (top-down). So food-web structure, not just nutrients, decides how far a
lake recovers.
\end{chain}

\textbf{Why recovery lags} after loading is cut: (i) the water must be exchanged; (ii) phosphorus stored in sediments keeps
being released, even faster as the water concentration falls; (iii) at peak eutrophication algae were limited by N or
light, so P must fall a long way before it limits; (iv) long-lived planktivorous fish keep suppressing grazers.

\subsection{Not so simple: alternative stable states}
\textbf{Lake Veluwe} (Netherlands): as phosphorus rose, charophyte vegetation collapsed abruptly (late 1960s--70s); when
loading was cut, the vegetation returned only in the 1990s, at a \textbf{much lower} phosphorus level.
\textbf{Hysteresis}: the path back differs from the path out.

\begin{chain}[: why a shallow lake has two stable states]
\begin{enumerate}
\item Sign rule (slide 16): multiply the signs round a loop. Turbidity $\to$ light on the bed ($-$) $\to$ vegetation
($+$) $\to$ sediment resuspension ($-$) $\to$ turbidity ($+$): $(-)(+)(-)(+)=+$. The direct pair vegetation
$\dashv$ turbidity $\dashv$ vegetation gives $(-)(-)=+$ too. The loop is \textbf{reinforcing}.
\item A reinforcing loop alone means runaway. But it saturates at both ends: plant cover cannot exceed the bed, and
clear water has no more turbidity to remove; once plants are gone, the plant part of the loop has nothing left to act on.
\item Reinforcing loop + saturation at both ends = \textbf{two self-stabilising states}: clear and vegetated, or turbid
and unvegetated. A small disturbance from either is pushed back by the same loop.
\item Plants keep water clear by taking up nutrients, sheltering \emph{Daphnia} from fish, and stopping resuspension.
Fish keep it turbid by eating \emph{Daphnia} and stirring sediment.
\item In \emph{deep} lakes light reaches the bed only near the shore, so plants cannot cover most of the lake and the loop is
too weak: deep lakes respond smoothly (background).
\end{enumerate}
\end{chain}

\textbf{Graphical model} (slide 17), three assumptions: turbidity rises with nutrients; vegetation lowers turbidity;
vegetation vanishes above a critical turbidity. Result: at low nutrients only the clear state exists, at high nutrients
only the turbid state, and in between \textbf{both}, separated by an unstable equilibrium. Crossing the edge of the
overlap gives a \textbf{catastrophic regime shift} (tipping point). Making the model more realistic does not change
this qualitative conclusion.

\begin{center}
\begin{tikzpicture}
\begin{axis}[width=8.5cm,height=5.6cm,xlabel={nutrient loading $a$},ylabel={turbidity $x$},
  xmin=-0.1,xmax=1.0,ymin=0,ymax=2.6,tick label style={font=\scriptsize},title={\small Box 1: $\dot x=a-x+1.4\,x^8/(x^8+1)$}]
\addplot[thick,NavyBlue,domain=0:0.723,samples=60,variable=\t] ({\t-1.4*\t^8/(\t^8+1)},{\t});
\addplot[thick,dashed,gray,domain=0.723:1.268,samples=60,variable=\t] ({\t-1.4*\t^8/(\t^8+1)},{\t});
\addplot[thick,BrickRed,domain=1.268:2.5,samples=60,variable=\t] ({\t-1.4*\t^8/(\t^8+1)},{\t});
\draw[-{Stealth},orange,thick] (axis cs:0.626,0.75) -- (axis cs:0.626,1.95);
\draw[-{Stealth},ForestGreen,thick] (axis cs:0.05,1.24) -- (axis cs:0.05,0.1);
\node[font=\scriptsize,anchor=west] at (axis cs:0.64,1.3) {collapse at 0.626};
\node[font=\scriptsize,anchor=west] at (axis cs:0.07,0.6) {recovery only at 0.050};
\end{axis}
\end{tikzpicture}

\small Solid blue: clear stable state; red: turbid stable state; dashed: unstable. Between the folds both states exist.
\end{center}

\begin{chain}[: hysteresis in numbers (Box 1 of Scheffer et al., background but explains the slide)]
Equilibria satisfy $bx-a=rf(x)$: a line of slope $b$ against an S-shaped recycling curve. Three crossings (two stable,
one unstable) are possible only if the S-curve is somewhere steeper than the line: $\max rf'(x)>b$. With $b=1$, $r=1.4$,
$h=1$, $p=8$ the bistable window is $0.050<a<0.626$. Raise loading past 0.626: the clear state disappears and turbidity
jumps $\sim$2.8-fold. Lower it again: the lake stays turbid until $a<0.050$. Restoring clarity by loading alone needs a
\textbf{$\sim$92\,\% cut} from the collapse value; returning to the pre-collapse loading does nothing.
\end{chain}

\begin{keyidea}[: resilience]
Phosphorus piling up in soil and mud need not change how the lake looks; it makes the clear state's ``valley'' shallower
(loss of \textbf{resilience}) until an ordinary shock (flood, hot summer, loss of piscivorous fish) tips it over. Such
shifts come ``quite unannounced''. The same mathematics (catastrophe theory, folds) applies to deserts, coral reefs,
markets, and the ice--albedo snowball of climate science. Near a fold, recovery from small shocks slows down
(critical slowing down), the one early warning theory predicts.
\end{keyidea}

\subsection{Restoration}
\textbf{Biomanipulation}: ``reduce the number of fishes''. Fewer planktivorous (and bottom-stirring) fish $\Rightarrow$
more \emph{Daphnia} $\Rightarrow$ phytoplankton grazed down $\Rightarrow$ clear water $\Rightarrow$ plants recolonise and
then hold the lake clear. A \emph{temporary} push can have a \emph{permanent} effect because it moves the lake into the other
basin of attraction, \textbf{provided the nutrient level is not too high} (otherwise the clear state does not exist).
Evidence: 18 Dutch lakes over ten years: transparency increased.

\begin{center}\small
\begin{tabularx}{\textwidth}{llXX}
\toprule
& method & principle & problem\\\midrule
physical & sediment dredging & remove P-rich mud (internal load) & expensive; disposal sites\\
& dilution/flushing & flush with P-poor water & needs large volumes of such water\\
& water-level manipulation & higher: less resuspension; lower: favours plants & not always easy or permitted\\
& hypolimnetic withdrawal & draw off P-rich bottom water & harms downstream waters\\
chemical & hypolimnetic oxygenation & raise redox so sediment holds P & $>$10 years; permanence hard\\
& P inactivation & Al, Fe, Ca salts bind P & Al toxic at low and high pH\\
biological & remove planktivorous fish & more large zooplankton, more grazing & permanence hard\\
& stock piscivores & predators cut planktivores & timing is crucial\\
& introduce mussels & filter the water (\emph{Dreissena}) & invasive, can explode\\\bottomrule
\end{tabularx}
\end{center}

\begin{numbers}
$P_{tot}$ classes: 5 / 10 / 30 / 100\,µg\,L$^{-1}$. Sewage N:P $\approx$ 4:1. Lake Washington: diversion 1968,
\emph{Daphnia} 1976, transparency 1 $\to$ 3--4 $\to$ 10\,m. Global P to fresh water $\approx$1.5\,Tg/yr (China 30\,\%,
India 8\,\%, USA 7\,\%, Spain and Brazil 6\,\% each; 2018 review). Biomanipulation: 18 Dutch lakes, 10 years.
Summary slide: fresh water threatened by bad policy (Aral Sea) and eutrophication; solutions need scientists,
politicians and, above all, local people.
\end{numbers}

\begin{statements}
\begin{enumerate}[label=\arabic*.]
\item ``Since phosphorus caused the turbid state, reducing phosphorus back to its pre-collapse level restores clear water.''
\F Hysteresis: recovery needs a much lower level (Veluwe; Box 1's $\sim$92\,\% cut).
\item ``A loop with an even number of negative links is reinforcing.''
\T The product of signs is positive.
\item ``A reinforcing feedback loop by itself guarantees two stable states.''
\F It gives runaway; two states need saturation (limits) at both ends.
\item ``Biomanipulation works by reducing nutrient loading.''
\F It changes the food web (fewer planktivores, more grazers) and pushes the lake across the threshold; nutrients
need not change, though they must not be too high.
\item ``Deep lakes show the same pronounced hysteresis as shallow lakes.''
\F Plants cannot cover a deep lake's bed, so the plant--turbidity loop is weak and response is smooth.
\item ``Recovery of Lake Washington tracked phosphorus alone.''
\F The second jump (to 10\,m) followed the \emph{Daphnia} return.
\item ``Hypolimnetic oxygenation suits a shallow, unstratified lake.''
\F It has no hypolimnion; dredging, P inactivation or biomanipulation fit.
\item ``A lake losing resilience will show visible deterioration before it flips.''
\F Shifts typically come unannounced; the state changes little before the fall.
\item ``Internal loading can keep a lake eutrophic after external loading stops.''
\T Sediment P keeps being released, especially under anoxia; it is one reason for the lag and the case for dredging.
\end{enumerate}
\end{statements}

\begin{longans}[: ``Reducing phosphate alone is often not enough to restore a lake. Explain.'']
(1) Eutrophication chain with internal loading as a reinforcing loop. (2) Lags: water renewal, sediment P, N/light
limitation, long-lived fish. (3) Alternative stable states in shallow lakes: sign rule, saturation, two states;
graphical model; hysteresis (Veluwe). (4) Lake Washington: second jump with \emph{Daphnia}; food-web control.
(5) Remedies: biomanipulation as a threshold-crossing push (works only if nutrients are not too high); dredging or
P inactivation to remove the internal stock; each method's drawback. (6) Resilience and unannounced shifts; need for
scientists, policy and local people.
\end{longans}
